%%
%% THESIS/DISSERTATION TEMPLATE FOR THE UNIVERSITY OF ALABAMA.
%%
%% This example is written by Paul Kilgo. It highlights all the (few) features
%% of uathesis.cls
%%

%% Use one of the following if writing a thesis/dissertation.
% \documentclass[thesis]{uathesis}
\documentclass[dissertation]{uathesis-2024}

%% Basic packages you'll probably want to use.
\usepackage{graphicx}                 %% For using \includegraphics{}
\usepackage{cite}                     %% For sorting/collapsing citations.
\usepackage{color}                    %% For colors used in listings.
\usepackage{listings}                 %% Code listings (for engineers)

%% Includes
\input{inc/glossary.tex}
% TODO: If you have any special commands defined, include them here.
\usepackage[titletoc]{appendix}%
\usepackage{kantlipsum}
%% Required parameters (these default to undefined)
\author{Student Name}       %% Your name!
\adviser{ Committee Member/}%Dissertation/}   %% Your adviser/committee chair!

%% The people on your committee.
%% Use \and to break them up between lines.
\committee{
  Committee Member 1/ Committee Co-Chair (if applicable) \and
  Committee Member 2 \and
  Committee Member 3 \and
  Committee Member 4
}

%% Note the use of \and to create line breaks and the
%% inverted pyramid style requested by the graduate school.
%% You have to do this manually. Sorry.
\title{Title of the thesis: the top line should be the longest \and
  the middle one is second longest \and
  and the last is shortest}

\degree{Master of Arts/Master of Science/Doctor of Philosophy/Doctor of Education}   %% Change to suit your degree.
\department{INSERT YOUR DEPARTMENT NAME HERE}   %% Change to suit your department.

%% These are body text paragraphs to be placed in the front matter.
\abstract{
  Type input{} your abstract here. 
  
  Abstract must not exceed 350 words. This page is always the first page of the preliminaries section. It is always the first numbered page, all preliminary pages are formatted with Roman numerals, bottom center of each page in sequence. This page is always ii.

  Set up the chapter/section header with single-line spacing, no indent, zero points before, 12 points after, with two inch top margin on the first page of the section.
  
  Where this section runs to more than one page those after page 1 have a one inch top margin.

 For the narrative return to regular page format{ }---{ }first line indent, double line spacing, 12 points after.

\textit{Keywords:} [Insert keywords here]
}

\dedication{
Type or input{} your dedication here. This section is optional, but if included it is always located immediately after the Abstract.
  
 Set up the chapter/section header with single line spacing, no indent, zero points before, 12 points after, with two-inch top margin on the first page of the section.

 Where this section runs to more than one page those after page 1 have a one inch top margin.

 For the narrative return to regular page format{ }---{ }first line indent, double line spacing, 12 points after.

}

\acknowledgments{
Type your acknowledgments here.

  This is always required. Include here any acknowledgements or affirmations required in association with any funding support for this research as well as any help or assistance provided by friends, faculty, and colleagues.
Set up the chapter/section header with single-line spacing, no indent, zero points before, 12 points after, with two-inch top margin on the first page of the section.
Where this section runs to more than one page those after page 1 have a one inch top margin.
For the narrative return to regular page format{ }---{ }first line indent, double line spacing, 12 points after.



}

%% Optional parameters. The default usage is shown.
\university{The University of Alabama}
\school{Graduate School}
\gradyear{\the\year}
\place{Tuscaloosa, Alabama}

\begin{document}

%% Makes the title, abstract, dedication, table of contents, etc.
%% Must be before \begin{body}
\makefrontmatter

%% BODY PORTION
%% The bulk of your thesis.
%% Begin your \chapter's here.

\begin{body}

%% Body chapters.
\chapter{INTRODUCTION}
Type or input{} introduction here. Note that the list of abbreviations and symbols is included when the abbreviations and symbols used in your manuscript are NOT self-explanatory. 


The Table of Contents is always required. 


Take a look at Figure \ref{fig:shapes}. It has some shapes. They're very
complex and interesting shapes. I could gaze at them for hours. Another interesting shape is shown in Figure \ref{fig:complex_shape} 

\begin{figure}[h]
  \centering
  \includegraphics[width=5.0in]{fig/shapes.pdf}
  \caption{Caption for the interesting shapes.}
  \label{fig:shapes}
\end{figure}

\begin{figure}[h]
  \centering
  \includegraphics[width=3.0in]{fig/rhombicosidodecahedron.png}
  \caption{A CAD rendering of a rhombicosadodecahedron. This caption is much longer in order to show how it looks in the List of Figures.}
  \label{fig:complex_shape}
\end{figure}

\chapter{LITERATURE REVIEW}

Type or input{} a chapter and include citations as needed \cite{rocket-themoon2167,lesk:1977}.
\section{Level 2 Heading}
Follow your chosen style guide for protocol for chapter titles (Numbered? Not numbered? Include/do not include word Chapter?). Section titles longer than one line must be single spaced. The headings in the text must match exactly with the headings in CONTENTS as to the page numbers, spelling, and punctuation.  However, for subheadings you may use bolding and italics in the text but you may not use bolding and italics in CONTENTS.

\subsection{Level 3 heading}
You can even add sub-subheadings.

\chapter{APPROACH (BECAUSE A LONGER-THAN-1-LINE CHAPTER TITLE IS IMPORTANT)}
Take a look at Table \ref{tab:love-shapes}. It shows how much I love certain
shapes. A second table is shown in \ref{tab:love-all}. This compares my love for shapes with my love for other things in life.

\begin{table}[h]
  \centering
  \input{tab/love-shapes.tex}
  \caption{Shapes and corresponding love in Love Standard Units (LSU).This table has a long caption that is important for highlighting the formatting.}
  \label{tab:love-shapes}
\end{table}

\begin{table}[h]
  \centering
  \input{tab/love-all.tex}
  \caption{Other objects and corresponding love in Love Standard Units (LSU).}
  \label{tab:love-all}
\end{table}

\chapter{EXPERIMENT}
Type or input{} a chapter. \cite{latex2e,knuth:1984}

\chapter{DISCUSSION}
Type or input{} a chapter. \cite{latex:companion}

\chapter{CONCLUSION}
Type or input{} a chapter. \cite{texbook}

%% Currently adding the references to the table of contents is not automatic
%% So we are forced to do it manually.

%% TODO: This is not a fix. This line has to come after the \chapter command
%% in order for the correct page to be inferred for the table of contents.
\clearpage
\addcontentsline{toc}{chapter}{REFERENCES}

 \begingroup
 \setstretch{1.0}
 \topskip=1in
\bibliographystyle{plain}
\bibliography{thesis-template-2024}

 \endgroup

%% Generates the bibliography.


\begin{appendices}
\addtocontents{toc}{\setlength{\cftchapnumwidth}{1em}}
\renewcommand{\appendixname}{APPENDIX}
\addtocontents{toc}{\setcounter{tocdepth}{0}}
\addtocontents{toc}{\protect\renewcommand{\protect\cftchappresnum}{}}

\chapter{INSERT TITLE}
Type or input{} first appendix here. Appendices must have a title, in all caps, and be ordered in some fashion.  For instance, APPENDIX A: SUPREME COURT CASES, APPENDIX B: APPEALS COURT CASES, APPENDIX C: TRIAL COURT CASES, etc. 

\section{IRB Approval Letters}
If your research has IRB approval letters they must be included as the final appendix (if you have other appendices) and it must be titled and formatted like the other appendices. The CONTENTS must signify what page has the IRB Approval Letters in it.
All ink signatures and initials are blacked out. Begin with a 2” top margin.

\end{appendices}

\end{body}
\end{document}